\documentclass[11pt]{article}
\usepackage[margin=1in]{geometry}
\usepackage{amsmath,amssymb,amsthm}
\usepackage[hidelinks]{hyperref}
\emergencystretch=2em
\newtheorem{theorem}{Theorem}
\newtheorem{lemma}{Lemma}
\newcommand{\Ord}{\mathrm{Ord}}
\newcommand{\Fix}{\mathrm{Fix}}
\newcommand{\nfp}{\mathrm{nfp}}
\title{Proof of conjecture 00000003527:\\The closed unbounded fixed-point chain of a normal function}
\author{GitHub: gaochengzhecpu}
\date{}
\begin{document}
\maketitle

\begin{abstract}
For every normal ordinal function, its fixed points form a closed unbounded
class and have a normal increasing enumeration. Its $\omega$-th enumerand is
the supremum of the finite enumerands. Separately, the least fixed point
above any starting ordinal is the supremum of the finite iterates of the
original function. We prove both statements and distinguish their two
sequences explicitly. These are classical facts; the Lean verification
uses the corresponding substantive Mathlib theorems.
\end{abstract}

\section*{Statement and the two limit constructions}
A function $f:\Ord\to\Ord$ is \emph{normal} when it is strictly increasing
and satisfies $f(\lambda)=\sup_{\alpha<\lambda}f(\alpha)$ at every nonzero
limit ordinal $\lambda$. Write $F=\{\beta:f(\beta)=\beta\}$.
For an ordinal $a$, let $\nfp_f(a)$ be the least member of $F$ at least $a$.
Its existence will be proved, rather than assumed. Let $d$ be the increasing
enumeration of $F$, indexed from zero.

\begin{theorem}\label{thm:main}
For every normal $f$, the class $F$ is closed in the order topology and
unbounded. Its enumeration $d$ is normal and has range exactly $F$, with
\begin{align*}
 d(0)&=\nfp_f(0),\\
 d(\alpha+1)&=\nfp_f(d(\alpha)+1),\\
 d(\lambda)&=\sup_{\alpha<\lambda}d(\alpha)
       &&(\lambda\text{ a nonzero limit}).
\end{align*}
In particular,
\begin{equation}\label{eq:enum}
 d(\omega)=\sup_{n<\omega}d(n).
\end{equation}
For every starting ordinal $a$, the other iteration formula is
\begin{equation}\label{eq:iterate}
 \nfp_f(a)=\sup_{n<\omega}f^{\circ n}(a).
\end{equation}
This is a fixed point at least $a$, and it is at most every fixed point
$b\geq a$.
\end{theorem}

Both source languages refer to the fixed-point sequence and its
$\omega$-th term. We read this sequence as the standard increasing
enumeration $d$, whose transfinite limit step is \eqref{eq:enum}.
Formula \eqref{eq:iterate} also addresses iteration of the original normal
function. These are distinct constructions: even for $f=\mathrm{id}$,
$d(\omega)=\omega$ whereas $\nfp_f(0)=0$. Thus no assertion identifies
the $\omega$-th fixed-point enumerand with iteration of $f$ from zero.

\section*{Proof}
First $\alpha\leq f(\alpha)$ for every ordinal $\alpha$. Transfinite
induction proves this: it is immediate at zero, strict increase gives the
successor step, and the supremum formula gives the limit step.
Normality also implies preservation of suprema of nonempty increasing
ordinal sequences. To see this, let $s$ be their supremum. If $s$ is
attained, the assertion follows by monotonicity. Otherwise $s$ is a limit
and the sequence is cofinal below $s$, so normality gives the assertion.

Fix $a$, put $x_0=a$ and $x_{n+1}=f(x_n)$, and set
$\beta=\sup_{n<\omega}x_n$. The sequence is increasing. Supremum
preservation and deletion of its first term give
\[
 f(\beta)=\sup_{n<\omega}f(x_n)
          =\sup_{n<\omega}x_{n+1}=\beta.
\]
Also $a\leq\beta$. If $b\geq a$ is fixed, induction gives $x_n\leq b$
for every $n$, so $\beta\leq b$. This proves both existence and leastness
of $\nfp_f(a)$ and proves \eqref{eq:iterate}. Applying it with starting
ordinal $a+1$ supplies a fixed point strictly larger than $a$. Hence $F$
is unbounded.

For closure, a normal function is continuous for the ordinal order topology.
At zero and successor ordinals the domain point is isolated. At a nonzero
limit $\lambda$, every threshold below $f(\lambda)$ is exceeded by
$f(\alpha)$ for some $\alpha<\lambda$; monotonicity on
$(\alpha,\lambda+1)$ then gives continuity at $\lambda$.
The ordinal order topology is Hausdorff, so the equalizer of the continuous
maps $f$ and the identity is closed. This equalizer is precisely $F$.
Equivalently, any nonempty increasing sequence of fixed points has a fixed
point as its supremum, by the same supremum-preservation calculation.

Define $d$ by the three transfinite recursion clauses of the theorem.
The initial and successor values are fixed by the preceding construction.
At a limit, the supremum of the preceding fixed points is fixed, so all
$d(\alpha)$ belong to $F$. Moreover
$d(\alpha+1)\geq d(\alpha)+1$. Transfinite induction, using the limit
clause, therefore proves strict increase; continuity at limits is a defining
clause. Thus $d$ is normal and in particular $\alpha\leq d(\alpha)$.

To check that this enumeration omits no fixed point, take $b\in F$.
There exists an index $\alpha$ with $b\leq d(\alpha)$, since $b\leq d(b)$.
Choose the least such index. If it is zero, leastness of
$\nfp_f(0)$ gives $d(0)\leq b$. If it is a successor $\gamma+1$,
minimality gives $d(\gamma)<b$, hence $d(\gamma)+1\leq b$, and the
least-fixed-point property gives $d(\gamma+1)\leq b$. If it is a
nonzero limit, all preceding values are below $b$, and their supremum is
at most $b$. In every case $d(\alpha)\leq b$, so equality holds.
Consequently the range of $d$ is exactly $F$.

Finally, apply the limit recursion to $\omega$. Its predecessors are
exactly the finite ordinals, giving \eqref{eq:enum}. This completes the
proof of Theorem~\ref{thm:main}.

\section*{Formal verification and scope}
The file \texttt{lean/Main.lean} uses Mathlib's actual universe-polymorphic
ordinals, normal functions, fixed-point set, ordinal derivative (the
fixed-point enumeration), and next-fixed-point construction. Closure is
the genuine topological predicate \texttt{IsClosed}; unboundedness is
the negation of \texttt{BddAbove}. An additional theorem supplies a fixed
point strictly beyond every ordinal.

The main theorem is
\begin{center}\texttt{Conjecture3527.conjecture\_3527\_true}.\end{center}
It proves the universally quantified \texttt{FixedPointClaim}: for every
normal ordinal function, topological closure, unboundedness, the exact
enumeration range, normality of the enumeration, its zero/successor/limit
clauses, its $\omega$-th value, and the finite-iteration formula with the
least-fixed-point property for every starting ordinal all hold.
The two values for the identity example are also formalized.

This is a formal verification of a classical theorem, not a claim to a new
fixed-point theorem. Substantial library results are used directly:
\begin{center}\small
\texttt{not\_bddAbove\_fp}, \texttt{IsNormal.fp\_iff\_deriv},\\
\texttt{isNormal\_deriv}, the recursion lemmas for \texttt{deriv},\\
\texttt{IsNormal.apply\_omega0}, \texttt{iSup\_iterate\_eq\_nfp},\\
\texttt{IsNormal.nfp\_fp}, and \texttt{nfp\_le\_fp}.
\end{center}
For closure, the proof combines the library equivalence between normality
and strict monotonicity plus continuity with the topological equalizer
theorem. All conclusions of the theorem are checked in Lean; there is no
additional paper-only premise. The paper gives the accompanying classical
argument rather than claiming its individual proof steps were reimplemented
independently of those library results.

\section*{Source and reproduction}
The exact English and Chinese source is preserved in \texttt{SOURCE.md},
from upstream commit
\begin{center}\texttt{b355e42906e6eaf52bc7e9d005e0ef02d8aa7400}.\end{center}
Its SHA-256 is
\begin{center}\small\texttt{1b36c16e1f37c47498803fa844b3550ed91df0b7920f33a0a7f852d47620c34a}.\end{center}
The parent source/history check found no earlier PR or solution.
Lean is 4.19.0 and Mathlib is pinned to
\begin{center}\texttt{c44e0c8ee63ca166450922a373c7409c5d26b00b}.\end{center}
The mathematical library sources are
\begin{center}\small
\path{SetTheory/Ordinal/FixedPoint.lean}\\
\path{SetTheory/Ordinal/Family.lean}\\
\path{SetTheory/Ordinal/Topology.lean}.
\end{center}
In the \texttt{lean} directory, run
\begin{center}\texttt{lake env lean -DwarningAsError=true Main.lean}\end{center}
and \texttt{lake build}. The author strict check passes; the audited
theorems use only \texttt{propext}, \texttt{Classical.choice}, and
\texttt{Quot.sound}. Compile this document with Tectonic or a compatible
LaTeX engine. Developed by a delegated AI agent with a self-review; the
coordinating agent reviews it separately. Review mode: agent development
with parent adversarial review; no external independent review.
\end{document}
